% Part 1 -- Medium Worked Problems

\begin{problem}[The Nested Quadratic Chain]
Consider the quadratic function
\[
f(x)=x^2-6x+7.
\]
Find all real values of $x$ satisfying
\[
f(f(x))=-1.
\]
\end{problem}

\begin{hint}
Do not expand $f(f(x))$ into a quartic straight away.
\begin{itemize}[leftmargin=*]
    \item Let $u=f(x)$ and solve $f(u)=-1$ first.
    \item Then set $f(x)=u$ for each value of $u$ you find.
    \item Completing the square is a clean way to solve the resulting quadratics.
\end{itemize}
\end{hint}

\begin{solution}
Let $u=f(x)$. Then $f(f(x))=-1$ becomes
\[
f(u)=-1,
\]
so
\[
u^2-6u+7=-1
\quad\Longrightarrow\quad
u^2-6u+8=0
\quad\Longrightarrow\quad
(u-4)(u-2)=0.
\]
Hence
\[
u=4 \quad \text{or} \quad u=2.
\]

So we solve the two equations:

\textbf{Case 1:} $f(x)=4$
\[
x^2-6x+7=4
\quad\Longrightarrow\quad
x^2-6x+3=0.
\]
Thus
\[
x=\frac{6\pm\sqrt{36-12}}{2}=3\pm\sqrt6.
\]

\textbf{Case 2:} $f(x)=2$
\[
x^2-6x+7=2
\quad\Longrightarrow\quad
x^2-6x+5=0
\quad\Longrightarrow\quad
(x-1)(x-5)=0.
\]
So
\[
x=1 \quad \text{or} \quad x=5.
\]

Therefore, the real solutions are
\[
\boxed{x=1,\ 5,\ 3-\sqrt6,\ 3+\sqrt6.}
\]
\end{solution}

\begin{takeaways}
    \begin{itemize}[leftmargin=*]
        \item \textbf{Level-by-level solving:} Although $f(f(x))$ would expand to a quartic, it is much better to treat it as a quadratic in disguise by introducing a substitution.
        \item \textbf{Composition structure:} For a composite equation $f(g(x))=c$, solve the outer layer first, then feed those answers back into the inner layer.
        \item \textbf{Root count intuition:} A quadratic composed with another quadratic can produce up to four real solutions, which is exactly what happens here.
        \item \textbf{Higher-ed connection:} Problems involving $f(f(x))$, $f(f(f(x)))$, and beyond are the first step toward iterated functions, dynamical systems, and fixed-point ideas.
    \end{itemize}
\end{takeaways}

\begin{problem}[Deriving Functional Rules Through Substitution]
Let $f$ be a function such that
\[
f:\mathbb{R}\setminus\{1\}\to\mathbb{R}.
\]
The function satisfies
\[
f\!\left(\frac{x+1}{x-1}\right)=\frac{x}{x+1}
\]
for all real numbers $x\ne \pm 1$.
\begin{enumerate}[label=(\roman*)]
\item Find the value of $x$ such that $\dfrac{x+1}{x-1}=2$. Hence calculate $f(2)$.
\item By using the substitution $u=\dfrac{x+1}{x-1}$, express $f(x)$ in terms of $x$ wherever the rule is determined.
\end{enumerate}
\end{problem}

\begin{hint}
For (i), solve the inner equation first, then substitute that value of $x$ into the right-hand side. For (ii), solve
\[
u=\frac{x+1}{x-1}
\]
for $x$ in terms of $u$, then substitute into $\dfrac{x}{x+1}$.
\end{hint}

\begin{solution}
\textbf{(i)} Solve
\[
\frac{x+1}{x-1}=2.
\]
Then
\[
x+1=2x-2
\quad\Longrightarrow\quad
x=3.
\]
Substitute $x=3$ into the original equation:
\[
f(2)=\frac{3}{3+1}=\boxed{\frac34}.
\]

\textbf{(ii)} Let
\[
u=\frac{x+1}{x-1}.
\]
Solving for $x$,
\[
u(x-1)=x+1
\quad\Longrightarrow\quad
x(u-1)=u+1
\quad\Longrightarrow\quad
x=\frac{u+1}{u-1}.
\]
Hence
\[
f(u)=\frac{\frac{u+1}{u-1}}{\frac{u+1}{u-1}+1}
=\frac{\frac{u+1}{u-1}}{\frac{2u}{u-1}}
=\frac{u+1}{2u}.
\]
Replacing $u$ by $x$ gives
\[
\boxed{f(x)=\frac{x+1}{2x}}
\]
for $x\ne 0,1$.

The value $f(0)$ is not determined by the given equation, because the input $0$ would require
\[
\frac{x+1}{x-1}=0
\quad\Longrightarrow\quad
x=-1,
\]
but $x=-1$ is excluded.
\end{solution}

\begin{takeaways}
\begin{itemize}[leftmargin=*]
\item The input to $f$ is the whole expression inside the brackets, not the parameter $x$ itself.
\item To recover a hidden function rule, invert the inner expression and substitute back into the output side.
\item Always check which inputs are actually covered by the substitution; here the equation determines $f(x)$ for $x\ne 0,1$, but not $f(0)$.
\end{itemize}
\end{takeaways}

\begin{problem}[Composites with a Parameter]
Given the functions
\[
f(x)=x^2-k \qquad (k\in\mathbb{R}), \qquad g(x)=2x+1,
\]
\begin{enumerate}[label=(\alph*)]
\item find simplified algebraic expressions for $f(g(x))$ and $g(f(x))$,
\item find the exact value of $k$ such that the graphs of $y=f(g(x))$ and $y=g(f(x))$ intersect at exactly one point,
\item for this value of $k$, find the coordinates of the intersection point.
\end{enumerate}
\end{problem}

\begin{hint}
First compute the two composites. Then set them equal and use the discriminant condition for a quadratic to have exactly one real solution.
\end{hint}

\begin{solution}
\textbf{(a)} Since $g(x)=2x+1$,
\[
f(g(x))=(2x+1)^2-k=4x^2+4x+1-k.
\]
Also
\[
g(f(x))=2(x^2-k)+1=2x^2-2k+1.
\]

\textbf{(b)} Intersection points satisfy
\[
f(g(x))=g(f(x)).
\]
So
\[
4x^2+4x+1-k=2x^2-2k+1
\]
which simplifies to
\[
2x^2+4x+k=0.
\]
For exactly one point of intersection, this quadratic must have exactly one real root:
\[
\Delta=4^2-4(2)(k)=16-8k=0.
\]
Hence
\[
\boxed{k=2.}
\]

\textbf{(c)} Substituting $k=2$,
\[
2x^2+4x+2=0
\quad\Longrightarrow\quad
x^2+2x+1=0
\quad\Longrightarrow\quad
(x+1)^2=0,
\]
so $x=-1$.

Then
\[
y=f(g(-1))=f(-1)=(-1)^2-2=-1.
\]
Therefore the intersection point is
\[
\boxed{(-1,-1).}
\]
\end{solution}

\begin{takeaways}
\begin{itemize}[leftmargin=*]
\item Composite functions are often compared by setting their outputs equal.
\item “Exactly one intersection” usually leads to a discriminant equal to $0$.
\end{itemize}
\end{takeaways}

\begin{problem}[A Relation That Splits into Two Functions]
Consider the relation
\[
y^2-4xy+5x^2=9.
\]
\begin{enumerate}[label=(\alph*)]
\item By treating the equation as a quadratic in $y$, use the quadratic formula to express $y$ in terms of $x$.
\item Hence explain why this relation does not represent a single function, but rather two separate function branches.
\item Determine the maximum domain of $x$ for which this relation is defined.
\item Find the exact points where the two branches meet.
\end{enumerate}
\end{problem}

\begin{hint}
Think of the equation as
\[
y^2+(-4x)y+(5x^2-9)=0.
\]
The two branches meet when the square-root part becomes zero.
\end{hint}

\begin{solution}
\textbf{(a)} Viewing the equation as a quadratic in $y$,
\[
y^2-4xy+(5x^2-9)=0.
\]
So
\[
y=\frac{4x\pm\sqrt{(-4x)^2-4(1)(5x^2-9)}}{2}.
\]
This gives
\[
y=\frac{4x\pm\sqrt{16x^2-20x^2+36}}{2}
=\frac{4x\pm\sqrt{36-4x^2}}{2}
=2x\pm\sqrt{9-x^2}.
\]

\textbf{(b)} Since
\[
y=2x+\sqrt{9-x^2}
\qquad \text{or} \qquad
y=2x-\sqrt{9-x^2},
\]
the relation gives two possible $y$-values for most allowed $x$-values. Therefore it is not a single function, but two branches:
\[
f_1(x)=2x+\sqrt{9-x^2},
\qquad
f_2(x)=2x-\sqrt{9-x^2}.
\]

\textbf{(c)} For the square root to be defined,
\[
9-x^2\ge 0
\quad\Longrightarrow\quad
x^2\le 9
\quad\Longrightarrow\quad
-3\le x\le 3.
\]
So the maximum domain is
\[
\boxed{[-3,3].}
\]

\textbf{(d)} The branches meet when
\[
\sqrt{9-x^2}=0,
\]
so
\[
x=\pm 3.
\]
Then $y=2x$, giving the points
\[
\boxed{(-3,-6)\quad \text{and}\quad (3,6).}
\]
\end{solution}

\begin{takeaways}
\begin{itemize}[leftmargin=*]
\item Some algebraic relations split naturally into two branches after solving for $y$.
\item A square root creates a domain restriction and often reveals where branches meet.
\end{itemize}
\end{takeaways}

\begin{problem}[A Family of Lines Through an Intersection Point]
Consider two distinct, non-parallel lines
\[
L_1: a_1x+b_1y+c_1=0,
\qquad
L_2: a_2x+b_2y+c_2=0,
\]
with unique intersection point $P(x_0,y_0)$.
\begin{enumerate}[label=(\alph*)]
\item Prove that for any real constant $k$, the equation $L_1+kL_2=0$ also passes through $P$.
\item Explain algebraically why $L_1+kL_2=0$ can represent every line through $P$, except for the line $L_2$ itself.
\item The lines $3x-4y+2=0$ and $x+2y-4=0$ intersect at a point $P$. Without finding the coordinates of $P$, find the equation of the line that passes through both $P$ and the origin.
\end{enumerate}
\end{problem}

\begin{hint}
At the intersection point $P$, both $L_1$ and $L_2$ equal $0$. In part (c), force the family line to pass through $(0,0)$.
\end{hint}

\begin{solution}
\textbf{(a)} Since $P$ lies on both lines,
\[
L_1(x_0,y_0)=0
\qquad \text{and} \qquad
L_2(x_0,y_0)=0.
\]
Therefore
\[
\bigl(L_1+kL_2\bigr)(x_0,y_0)=L_1(x_0,y_0)+kL_2(x_0,y_0)=0+k\cdot 0=0.
\]
So $L_1+kL_2=0$ passes through $P$ for every real $k$.

\textbf{(b)} Any line through $P$ in this family can be written more generally as
\[
\lambda L_1+\mu L_2=0,
\]
where $\lambda$ and $\mu$ are not both zero. If $\lambda\ne 0$, divide through by $\lambda$ to obtain
\[
L_1+\left(\frac{\mu}{\lambda}\right)L_2=0.
\]
Thus every such line can be written in the form $L_1+kL_2=0$, except the special case $\lambda=0$, which gives precisely the line
\[
L_2=0.
\]

\textbf{(c)} Write the required line as
\[
(3x-4y+2)+k(x+2y-4)=0.
\]
Since it passes through the origin,
\[
2+k(-4)=0
\quad\Longrightarrow\quad
2-4k=0
\quad\Longrightarrow\quad
k=\frac12.
\]
Hence
\[
(3x-4y+2)+\frac12(x+2y-4)=0.
\]
Multiplying by $2$ gives
\[
6x-8y+4+x+2y-4=0,
\]
so the line is
\[
\boxed{7x-6y=0.}
\]
\end{solution}

\begin{takeaways}
\begin{itemize}[leftmargin=*]
\item The linear combination $L_1+kL_2=0$ is a standard way to describe a whole pencil of lines through a common point.
\item This method often avoids solving for the actual intersection point.
\end{itemize}
\end{takeaways}

\begin{problem}[Why a Quadratic Cannot Have Three Distinct Roots]
Show that the quadratic equation
\[
ax^2+bx+c=0
\qquad (a\ne 0)
\]
cannot have more than two distinct roots.

\textit{Hint provided in the question: suppose the equation has three distinct roots $\alpha,\beta,\gamma$, then substitute them into the equation and derive a contradiction. Also explain the result visually using the graph of a quadratic function.}
\end{problem}

\begin{hint}
Write down the three equations obtained from $\alpha,\beta,\gamma$. Subtract pairs of equations to eliminate $c$.
\end{hint}

\begin{solution}
Suppose, for contradiction, that
\[
ax^2+bx+c=0
\]
has three distinct roots $\alpha,\beta,\gamma$.

Then
\[
a\alpha^2+b\alpha+c=0,
\qquad
a\beta^2+b\beta+c=0,
\qquad
a\gamma^2+b\gamma+c=0.
\]
Subtract the first two equations:
\[
a(\alpha^2-\beta^2)+b(\alpha-\beta)=0.
\]
Since $\alpha\ne\beta$,
\[
(\alpha-\beta)\bigl(a(\alpha+\beta)+b\bigr)=0
\quad\Longrightarrow\quad
a(\alpha+\beta)+b=0.
\]
Similarly,
\[
a(\beta+\gamma)+b=0.
\]
Subtracting these gives
\[
a(\alpha-\gamma)=0.
\]
But $\alpha\ne\gamma$, so this forces
\[
a=0,
\]
which contradicts the assumption $a\ne 0$.

Therefore a quadratic cannot have three distinct roots.

\textbf{Visual explanation.} The graph of $y=ax^2+bx+c$ is a parabola. A parabola can intersect the $x$-axis at at most two points, so the equation $ax^2+bx+c=0$ can have at most two real roots.
\end{solution}

\begin{takeaways}
\begin{itemize}[leftmargin=*]
\item A degree-$2$ polynomial has at most two distinct roots.
\item The algebra matches the geometry: a parabola meets the $x$-axis at most twice.
\end{itemize}
\end{takeaways}

\begin{problem}[Symmetry About the Vertex]
Consider the quadratic function
\[
f(x)=x^2-2x-8.
\]
\begin{enumerate}[label=(\alph*)]
\item By completing the square, express $f(x)$ in vertex form and state the coordinates of the vertex.
\item Let $x=v$ be the $x$-coordinate of the vertex. Simplify $f(v+h)$ and $f(v-h)$.
\item What does the result prove about the geometric shape of the parabola relative to its vertex?
\end{enumerate}
\end{problem}

\begin{hint}
Complete the square first. Then use the vertex $x$-coordinate $v$ as the centre of the two shifted inputs.
\end{hint}

\begin{solution}
\textbf{(a)}
\[
f(x)=x^2-2x-8=(x-1)^2-9.
\]
So the vertex is
\[
\boxed{(1,-9).}
\]

\textbf{(b)} Here $v=1$. Then
\[
f(1+h)=(1+h)^2-2(1+h)-8=h^2-9,
\]
and
\[
f(1-h)=(1-h)^2-2(1-h)-8=h^2-9.
\]
Hence
\[
\boxed{f(v+h)=f(v-h).}
\]

\textbf{(c)} Since points equally spaced to the left and right of $x=v$ give the same $y$-value, the parabola is symmetric about the vertical line
\[
\boxed{x=v=x_{\text{vertex}}.}
\]
So the axis of symmetry is $x=1$.
\end{solution}

\begin{takeaways}
\begin{itemize}[leftmargin=*]
\item Completing the square reveals the vertex and axis of symmetry immediately.
\item The identity $f(v+h)=f(v-h)$ is a precise algebraic expression of mirror symmetry for a quadratic.
\item \textbf{The Cubic Connection:} A strikingly similar observation is seen with cubic functions. While a quadratic has line symmetry perfectly balanced around its vertex, a cubic polynomial has point symmetry (rotational symmetry) balanced perfectly around its point of inflection. If $\mu$ is the $x$-coordinate of a cubic's inflection point, stepping $h$ units left and right reveals the identity $f(\mu+h)+f(\mu-h)=2f(\mu)$. The algebraic logic of exploring symmetry through $+h$ and $-h$ shifts remains exactly the same.
\end{itemize}
\end{takeaways}

\begin{problem}[A Family of Quadratics and the Locus of the Vertex]
Consider the family
\[
y=x^2-2kx+(k^2+k-2),
\]
where $k$ is a real parameter.
\begin{enumerate}[label=(\alph*)]
\item Express the function in vertex form.
\item Hence write down the vertex in terms of $k$.
\item Find the Cartesian equation of the locus of the vertex as $k$ varies, and describe it geometrically.
\item Find the $x$-intercepts in terms of $k$, and state the restriction on $k$ for which they exist.
\item Find the exact range of $k$ for which the parabola has one strictly positive and one strictly negative $x$-intercept.
\end{enumerate}
\end{problem}

\begin{hint}
Complete the square in $x$. For part (e), think about what the signs of the two roots tell you about their product.
\end{hint}

\begin{solution}
\textbf{(a)}
\[
y=x^2-2kx+k^2+k-2=(x-k)^2+k-2.
\]

\textbf{(b)} Therefore the vertex is
\[
\boxed{(k,\ k-2).}
\]

\textbf{(c)} Let the vertex be $(x,y)$. Since $x=k$, we have
\[
y=k-2=x-2.
\]
So the locus is
\[
\boxed{y=x-2,}
\]
which is a straight line.

\textbf{(d)} Put $y=0$:
\[
x^2-2kx+(k^2+k-2)=0.
\]
Using the completed-square form,
\[
(x-k)^2=2-k.
\]
Hence
\[
x=k\pm \sqrt{2-k}.
\]
These are real only when
\[
2-k\ge 0
\quad\Longrightarrow\quad
\boxed{k\le 2.}
\]

\textbf{(e)} For one root to be positive and the other negative, their product must be negative:
\[
k^2+k-2<0.
\]
Factorising,
\[
(k+2)(k-1)<0.
\]
Therefore
\[
\boxed{-2<k<1.}
\]
\end{solution}

\begin{takeaways}
\begin{itemize}[leftmargin=*]
\item A parameter family often becomes much clearer after completing the square.
\item Tracking the vertex is a standard route to finding a locus.
\item Opposite-sign roots correspond to a negative product.
\end{itemize}
\end{takeaways}

\begin{problem}[Composition of One-to-One Functions]
Prove or disprove the statement:

\begin{center}
\emph{The composition of two one-to-one functions is always another one-to-one function.}
\end{center}
\end{problem}

\begin{hint}
To prove a function is one-to-one, start from
\[
(f\circ g)(x_1)=(f\circ g)(x_2)
\]
and use the fact that both $f$ and $g$ are one-to-one.
\end{hint}

\begin{solution}
The statement is \textbf{true}.

Let $f$ and $g$ be one-to-one functions. We want to show that
\[
f\circ g
\]
is also one-to-one.

Assume
\[
(f\circ g)(x_1)=(f\circ g)(x_2).
\]
Then
\[
f(g(x_1))=f(g(x_2)).
\]
Since $f$ is one-to-one, equal outputs imply equal inputs. Therefore
\[
g(x_1)=g(x_2).
\]
Now use the fact that $g$ is one-to-one. Again, equal outputs imply equal inputs, so
\[
x_1=x_2.
\]
Hence
\[
(f\circ g)(x_1)=(f\circ g)(x_2)\implies x_1=x_2,
\]
which proves that $f\circ g$ is one-to-one.

Therefore, the composition of two one-to-one functions is always one-to-one.
\end{solution}

\begin{takeaways}
\begin{itemize}[leftmargin=*]
\item To prove a composition is one-to-one, peel the functions off from the outside in.
\item Injective functions preserve distinctions: different inputs cannot collapse to the same output.
\item This result helps explain why inverse functions behave well under composition.
\end{itemize}
\end{takeaways}

\begin{problem}[Domain Overlaps and One-to-One Testing]
Consider the following two relations, defined for all real numbers $x$:

\[
S(x)=
\begin{cases}
2x & \text{for } x \text{ rational},\\
x^3 & \text{for } x \text{ irrational},
\end{cases}
\]

\[
T(x)=
\begin{cases}
x & \text{for } x^2 \text{ rational},\\
-x & \text{for } x \text{ irrational}.
\end{cases}
\]

\begin{enumerate}[label=(\alph*)]
\item Determine whether $S(x)$ and $T(x)$ are valid functions. Give reasons.
\item For any relation from part (a) that is a valid function, decide whether it is one-to-one or many-to-one.
\item For any valid function, find all values of $x$ such that $f(x)=x$.
\end{enumerate}
\end{problem}

\begin{hint}
For a valid function, every input must produce exactly one output. In $T(x)$, check whether the two conditions can overlap. In part (b), try to find two different inputs for $S(x)$ that both produce the same output.
\end{hint}

\begin{solution}
\textbf{(a)} Since every real number is either rational or irrational, but not both, the two cases of $S(x)$ are disjoint and cover all real numbers. Hence
\[
\boxed{S(x)\text{ is a valid function}.}
\]

For $T(x)$, take $x=\sqrt2$. Then $x$ is irrational, so one rule gives
\[
T(\sqrt2)=-\sqrt2.
\]
But $x^2=2$ is rational, so the other rule gives
\[
T(\sqrt2)=\sqrt2.
\]
One input gives two outputs, so
\[
\boxed{T(x)\text{ is not a valid function}.}
\]

\textbf{(b)} To classify $S(x)$, note that
\[
S(1)=2
\]
and, since $\sqrt[3]{2}$ is irrational,
\[
S(\sqrt[3]{2})=(\sqrt[3]{2})^3=2.
\]
Thus
\[
S(1)=S(\sqrt[3]{2}) \quad \text{with} \quad 1\ne \sqrt[3]{2},
\]
so $S(x)$ is not one-to-one. Therefore
\[
\boxed{S(x)\text{ is many-to-one}.}
\]

\textbf{(c)} Solve $S(x)=x$ branch by branch.

If $x$ is rational, then
\[
2x=x \implies x=0,
\]
and this is valid since $0$ is rational.

If $x$ is irrational, then
\[
x^3=x \implies x(x^2-1)=0 \implies x=0,\pm1.
\]
But these are all rational, so none is allowed in the irrational branch.

Hence the only fixed point is
\[
\boxed{x=0.}
\]
\end{solution}

\begin{takeaways}
\begin{itemize}[leftmargin=*]
\item In a piecewise definition, the cases must give exactly one output for each input.
\item It is not enough for the rules to look different; you must check whether their conditions overlap.
\item When solving within a branch of a piecewise function, always test whether the solution really belongs to that branch.
\end{itemize}
\end{takeaways}
